% Release identifiers remain pending author decisions; example DOIs are not reused.
\section{Data and code availability}\label{sec:availability}

\paragraph{Research artifacts.}
The public author-review release provides executable source, saved result JSON, manuscript sources, generated tables and figures, and both PDFs at \url{https://github.com/cyberdid/archive-index-public}, tag \code{paper-author-review-2026-10-23}. No software/data or preprint DOI has been assigned. The generated-table manifest binds numerical artifacts to their source hashes. The public package supports rebuilding tables and PDFs from saved results. The full frozen corpus, query/run/assessment inputs and their 14-stage replay are retained in a separate research archive; they are not redistributed in this release while third-party attribution and scope review remain pending. Exact end-to-end statistical replay therefore requires those separately held inputs. The corpus-construction and evaluation scripts are public.

\paragraph{Reproduction.}
The isolated \path{paper/run_all.sh} runner distinguishes a saved-result rendering pass from a full frozen-input statistical replay. The completed full replay has 14 stages and reproduces ten reports, canonical qrels and recorded ANN input hashes. Physical cost, multimodal, ANN and private-field measurements are imported artifacts in this replay, not newly repeated experiments. The recorded 315 implementation tests, including a locally cached real-model integration check, address implementation behaviour separately. Observed environment specifications and hash-locked requirements accompany the code; a fresh full CUDA installation is not claimed. Rebuilding tables or PDFs does not by itself repeat model grading or physical measurements.

\paragraph{Source policy and attribution.}
The corpus manifest records source URL, collected licence label and SHA-256 for every file. Notes and rendered documents derive from Ukrainian and Russian Wikipedia (CC~BY-SA~4.0); the 150 selected data.gov.ua tables carry CC~BY labels. Crossmodal-3600 images retain per-image CC~BY~2.0 attribution, with CC~BY~4.0 captions; FLEURS recordings and transcripts carry CC~BY~4.0. Repository selection accepted MIT or Apache-2.0 API labels; the collected code repositories are labelled MIT in the manifest. That repository-level label is not a per-asset licence determination. Embedded third-party notices remain applicable: the attribution inventory records, among other cases, the Rubik font's SIL~OFL notice, a pinned missing-license supplement, and an application-terms scope question. Original notices, attribution files and the audit in \path{paper/corpus/LICENSES.md} are preserved. The final redistribution review of third-party material remains pending.

\paragraph{Licensing.}
The original executable source is licensed under Apache License~2.0. Article text, original figures and experimental records are licensed under Creative Commons Attribution~4.0 International (\url{https://creativecommons.org/licenses/by/4.0/}), to the extent the author holds the relevant rights. The licensing guide distinguishes these materials from third-party corpus assets and dependencies, which retain their own terms.

\paragraph{Private data and release inspection.}
Only aggregate operational checks are exported from the author's archive; its documents, query texts, paths, snippets and index are excluded. The release guard checks filenames, known marker bytes, decoded JSON, extracted PDF text/metadata and supported compressed content. Four investigated public word/name collisions require explicit source evidence and exact file hashes; full private-query and home-path markers have no exception. These checks do not constitute a universal privacy guarantee or OCR inspection of images. The existing private Git history is not an appropriate public release; publication requires a separately reviewed export.

\paragraph{Models and software.}
The experiment records identify EmbeddingGemma~2 (\code{google/embeddinggemma-2}, revision \code{914f7f8}), multilingual-E5-large (\code{intfloat/multilingual-e5-large}), the llama.cpp/GGUF builds and every judging model. Package, model and command fingerprints remain tied to their recorded runs. Model weights are not included in the research export.

\paragraph{AI-assisted work.}
Model-written queries and model relevance assessments are described in Section~\ref{sec:uapar} and the separate extension analysis; GPT grades remain comparison-only. Claude and Codex assisted implementation and manuscript preparation, and the Scientific Agent Skills workflow~\citep{kassis2026skills} informed source and claim checks. The author is responsible for the study design, interpretations, references and final declarations. No AI system is an author. This edition includes no human text-relevance audit; source-dataset multimodal pairings and implementation checks do not replace one.
